% Canonical Theta Field Definition
% Version: 2.1 - Phase 2 (Per COPILOT_INSTRUCTIONS)
% Date: 2025-11-14
% Status: Canonical - DO NOT DUPLICATE

\section{The Fundamental Biquaternion Field $\Theta(q,\tau)$}
\label{sec:canonical:theta_field}

\subsection{Definition}

The fundamental field of Unified Biquaternion Theory is a \textbf{biquaternion}:

\begin{equation}
\label{eq:canonical:theta_field}
\boxed{\Theta(q,\tau) \in \mathcal{B} = \mathbb{H} \otimes \mathbb{C}}
\end{equation}

\noindent where:
\begin{itemize}
    \item $q \in \mathcal{B}$ is a biquaternion coordinate with 4 degrees of freedom
    \item $\tau = t + i\psi$ is the canonical complex time
    \item Biquaternion time $T_B = t + i\psi + j\chi + k\xi$ is a deprecated/historical extension (see \texttt{speculative\_extensions/})
    \item $\mathcal{B} = \mathbb{H} \otimes \mathbb{C}$ denotes the biquaternion algebra
    \item $\mathbb{H}$ = quaternions with units $\{1, i, j, k\}$
\end{itemize}

\subsubsection{Explicit Biquaternion Form}

The field $\Theta$ is written explicitly as:

\begin{equation}
\label{eq:canonical:theta_explicit}
\Theta(q,\tau) = \Theta_0(q,\tau) + \Theta_1(q,\tau) i + \Theta_2(q,\tau) j + \Theta_3(q,\tau) k
\end{equation}

where $\Theta_a(q,\tau) \in \mathbb{C}$ ($a = 0,1,2,3$) are \textbf{complex-valued} component fields.

\textbf{Important}: Θ is a biquaternion (NOT a matrix). Matrix representations are computational tools only.

\subsection{Faithful Matrix Representation and Degree Count}

A biquaternion has four complex coefficients and therefore eight real
components:
\begin{equation}
  \Theta=\Theta_0+\Theta_1\mathbf i+\Theta_2\mathbf j+\Theta_3\mathbf k,
  \qquad \Theta_a\in\mathbb C.
\end{equation}
A faithful computational representation is a constrained $2\times2$ complex
matrix,
\begin{equation}
\label{eq:canonical:theta_matrix}
  \rho(\Theta)=
  \begin{pmatrix}
    \Theta_0+i\Theta_1 & \Theta_2+i\Theta_3\\
    -\Theta_2+i\Theta_3 & \Theta_0-i\Theta_1
  \end{pmatrix}.
\end{equation}
Thus
\begin{equation}
  \dim_{\mathbb C}\mathbb B=4,
  \qquad
  \dim_{\mathbb R}\mathbb B=8.
\end{equation}
A generic $4\times4$ or $8\times8$ complex matrix is not a biquaternion.
Matrix-valued flavour or Standard-Model extensions, when used, must be denoted
separately (for example $\Theta_{\rm ext}\in\operatorname{Mat}_n(\mathbb B)$)
and must not be silently substituted for the canonical core field.

\subsection{Hermitian Involution}

For variational and metric formulas we use the full Hermitian involution
$\ddagger$, which combines complex conjugation with quaternion conjugation and
corresponds to matrix Hermitian conjugation in the representation
\eqref{eq:canonical:theta_matrix}:
\begin{equation}
  \rho(\Theta^\ddagger)=\rho(\Theta)^\dagger.
\end{equation}
The notation must distinguish this operation from quaternion conjugation without
complex conjugation.

\subsection{Physical Interpretation}

The field $\Theta(q,\tau)$ encodes:

\begin{enumerate}
    \item \textbf{Geometric structure}: Biquaternionic tetrads $E_\mu$ emerge from $\Theta$ configurations, giving the raw quadratic tensor $\mathcal Q_{\mu\nu}=E_\mu E_\nu^\ddagger$; the observable metric is the normalized $\psi$-averaged real scalar Gram projection
    
    \item \textbf{Matter content}: Fermion fields and masses arise from internal structure
    
    \item \textbf{Gauge fields}: Electromagnetic and weak/strong interactions emerge from phase structure
    
    \item \textbf{Internal auxiliary sector}: Imaginary time component $\psi$ provides dynamics for an internal auxiliary field (speculative interpretations in \texttt{speculative\_extensions/})
\end{enumerate}

\subsection{Field Equations}

The field $\Theta$ satisfies the fundamental field equation (the \textbf{T-shirt formula}):

\begin{equation}
\label{eq:canonical:theta_equation}
\nabla^\dagger \nabla \Theta(q,\tau) = \kappa \mathcal{T}(q,\tau)
\end{equation}

\noindent where:
\begin{itemize}
    \item $\nabla$ is the \textbf{full covariant derivative} including both biquaternionic gravitational connection $\Omega_\mu$ (fundamental) and Standard Model gauge connections (see detailed structure below)
    \item $\nabla^\dagger \nabla$ is the gauge-gravity unified d'Alembertian operator in curved spacetime
    \item $\kappa$ is a coupling constant proportional to $8\pi G$ (Newton's constant)
    \item $\mathcal{T}(q,\tau)$ is the biquaternionic stress-energy source
\end{itemize}

\subsubsection{Structure of the Covariant Derivative}

The covariant derivative $\nabla_\mu$ acting on $\Theta$ is \textbf{not} an ordinary partial derivative. It is defined as:

\begin{equation}
\label{eq:canonical:nabla_structure}
\nabla_\mu \Theta = \partial_\mu \Theta + \Omega_\mu^{\mathrm{grav}} \Theta + A_\mu^{\mathrm{SM}} \Theta
\end{equation}

where:

\begin{enumerate}
    \item \textbf{Biquaternionic gravitational connection} $\Omega_\mu^{\mathrm{grav}} \in \mathbb{B}$ encodes fundamental spacetime curvature (Christoffel symbols/Levi-Civita connection arise as $\Gamma^\lambda_{\mu\nu} := \text{Re}(\Omega^\lambda_{\mu\nu})$ in the GR limit)
    
    \item \textbf{Standard Model gauge connection}:
    \begin{equation}
    A_\mu^{\mathrm{SM}} = i g_1 B_\mu Y + i g_2 W_\mu^a T^a + i g_3 G_\mu^A \Lambda^A
    \end{equation}
    
    containing $U(1)_Y$ hypercharge, $SU(2)_L$ weak, and $SU(3)_c$ strong interactions.
\end{enumerate}

Thus, the T-shirt equation \eqref{eq:canonical:theta_equation} is a \textbf{unified gauge-gravity field equation} where all fundamental interactions (gravity + electroweak + strong) are encoded within the single operator $\nabla$.

The operator $\nabla^\dagger\nabla$ in curved spacetime is precisely the Laplace-Beltrami/d'Alembertian operator that depends on both the metric (curvature) and the gauge field strengths. It includes:
\begin{itemize}
    \item Pure gravitational terms (Riemann curvature)
    \item Pure gauge terms (Yang-Mills field strengths for $U(1)$, $SU(2)$, $SU(3)$)
    \item Mixed gauge-gravity couplings
    \item Covariant derivatives of $\Theta$
\end{itemize}

For the complete explicit derivation of $\nabla$ structure, see Appendix on the structure of the covariant derivative.

\subsection{Normalization}

The field satisfies the normalization condition:

\begin{equation}
\label{eq:canonical:theta_normalization}
\text{Tr}(\Theta^\dagger \Theta) = \text{const.}
\end{equation}

for bound states. The constant depends on the physical system under consideration.

\subsection{Gauge Transformations}

Under local gauge transformations:

\begin{equation}
\label{eq:canonical:theta_gauge}
\Theta(q,\tau) \to U(q,\tau) \Theta(q,\tau) V^\dagger(q,\tau)
\end{equation}

\noindent where $U,V \in U(n)$ are unitary matrices encoding:
\begin{itemize}
    \item $U(1)$ electromagnetic gauge symmetry
    \item $SU(2)$ weak isospin symmetry
    \item $SU(3)$ color symmetry
\end{itemize}

\subsection{Reality Conditions}

For physical observables, we impose:

\begin{equation}
\label{eq:canonical:theta_reality}
\text{Re}\,\text{Tr}(\Theta^\dagger \mathcal{O} \Theta) \in \mathbb{R}
\end{equation}

for any observable operator $\mathcal{O}$.

\subsection{Asymptotic Behavior}

At spatial infinity ($|q| \to \infty$):

\begin{equation}
\label{eq:canonical:theta_asymptotics}
\Theta(q,\tau) \to \Theta_0 + O(1/|q|)
\end{equation}

\noindent where $\Theta_0$ is the vacuum configuration.

For temporal asymptotics (large $|t|$), boundary conditions depend on the specific physical scenario (scattering, bound states, etc.).

\subsection{Connection to Biquaternions}

The field can be expressed in biquaternion basis:

\begin{equation}
\label{eq:canonical:theta_biquaternion}
\Theta = \sum_{A=0}^{3} \sum_{B=0}^{3} \theta_{AB}(q,\tau) \, \sigma_A \otimes \sigma_B
\end{equation}

\noindent where $\sigma_A$ are Pauli matrices (with $\sigma_0 = I$) and $\theta_{AB}$ are complex scalar fields.

This representation makes the biquaternionic structure explicit.

\subsection{Units and Dimensions}

In natural units ($\hbar = c = 1$):

\begin{equation}
[\Theta] = [\text{mass}]^1 = [\text{length}]^{-1}
\end{equation}

This ensures dimensional consistency with standard field theory.

\subsection{Historical Note}

This definition supersedes all previous versions found in the repository, including:
\begin{itemize}
    \item 4D biquaternion representation (old preprint)
    \item Alternative spinor formulations
    \item Conflicting matrix dimensions
\end{itemize}

\textbf{This is the canonical version.} All other formulations should be considered historical or deprecated.

% End of canonical theta field definition
